% section: 分数一次変換と離散リカッチ方程式

\section{分数一次変換と離散リカッチ方程式}

\subsection{一次分数型は関数の反復である}

\[
  a_{n+1}=\frac{pa_n+q}{ra_n+s}
\]
を考える.
これは一次分数変換,またはメビウス変換と呼ばれる.
分母があるため一見扱いにくいが,関数の反復として見ると構造が現れる.

不動点は,
\[
  \alpha=\frac{p\alpha+q}{r\alpha+s}
\]
より,
\[
  r\alpha^2+(s-p)\alpha-q=0
\]
を満たす.
従って,不動点は二次方程式の解として求められる.

\subsection{二つの不動点を基準にする}

異なる二つの不動点を $\alpha,\beta$ とする.
次の変数変換を考える.
\[
  b_n=\frac{a_n-\alpha}{a_n-\beta}
\]
これは,$a_n$ が $\alpha$ と $\beta$ のどちらに近いかを,比として記録する量である.

変換
\[
  f(x)=\frac{px+q}{rx+s}
\]
について,固定点の条件を使うと,
\[
  f(x)-\alpha
  =
  \frac{(p-r\alpha)(x-\alpha)}{rx+s}
\]
および
\[
  f(x)-\beta
  =
  \frac{(p-r\beta)(x-\beta)}{rx+s}
\]
となる.
従って,
\[
  \begin{aligned}
    \frac{f(x)-\alpha}{f(x)-\beta}
     & =
    \frac{p-r\alpha}{p-r\beta}
    \frac{x-\alpha}{x-\beta}
  \end{aligned}
\]

ここで
\[
  \lambda=\frac{p-r\alpha}{p-r\beta}
\]
と置けば,
\[
  b_{n+1}=\lambda b_n
\]
となる.

分数を含む複雑な漸化式が,変数変換によって等比数列になった.
本編で扱った一次分数型の解法は,このような「二つの基準点からの比を取る」という幾何学的な意味を持っている.

\subsection{連続リカッチ方程式との共通点}

連続的なリカッチ方程式では,既知の特別解を基準にして,
\[
  x=x_0+\frac1u
\]
と置いた.
離散的な一次分数型では,二つの不動点を基準にして,
\[
  b=\frac{x-\alpha}{x-\beta}
\]
と置いた.

変換の形は違うが,発想は同じである.
\[
  \text{非線形で扱いにくい変化}
  \longrightarrow
  \text{基準となる解や点との差}
  \longrightarrow
  \text{線形または等比的な変化}
\]

漸化式と微分方程式を結ぶのは,公式の対応ではない.
「どの量を記録すれば,変化が単純になるか」という問題意識の対応である.
